% BHU PYQ -- AIHC & Archaeology: Mint and Minting of Coins in Ancient India (2025-26 NEP)
\documentclass[11pt,a4paper]{article}
\usepackage[a4paper,top=2.5cm,bottom=2.5cm,left=2.8cm,right=2.8cm,headheight=14pt]{geometry}
\usepackage{amsmath,amssymb}
\usepackage{fontspec}
\setmainfont{Kohinoor Devanagari}
\usepackage{microtype,fancyhdr,enumitem,hyperref,lastpage,parskip}

\hypersetup{colorlinks=true,urlcolor=black,linkcolor=black,pdftitle={AIHSE31: Mint and Minting of Coins in Ancient India -- BHU NEP 2025-26}}

\pagestyle{fancy}\fancyhf{}
\renewcommand{\headrulewidth}{0.4pt}\renewcommand{\footrulewidth}{0.4pt}
\fancyhead[L]{\small \textit{Banaras Hindu University}}
\fancyhead[R]{\small \textit{AIHSE31: Minting of Coins (2025-26)}}
\fancyfoot[C]{\small Page \thepage\ of \pageref{LastPage}}

\newlist{parts}{enumerate}{2}
\setlist[parts,1]{label=(\alph*),leftmargin=*,itemsep=4pt,topsep=3pt}
\setlist[parts,2]{label=(\roman*),leftmargin=*,itemsep=2pt,topsep=2pt}

\newcommand{\pts}[1]{\hfill \textit{[#1]}}

\begin{document}

\begin{center}
  {\large\bfseries BANARAS HINDU UNIVERSITY}\\[2pt]
  {\normalsize Under Graduate Semester III Examination, 2025-26 (NEP)}\\[6pt]
  \rule{0.8\linewidth}{0.5pt}\\[6pt]
  {\large\bfseries ANCIENT INDIAN HISTORY, CULTURE \& ARCHAEOLOGY}\\[4pt]
  {\large\bfseries Paper Code: AIHSE31 : Mint and Minting of Coins in Ancient India}\\[4pt]
  {\normalsize Time: 2:30 Hours \hfill Full Marks: 70}\\[2pt]
  {\small REF NO. 2025-26/D-III/071}
\end{center}
\rule{\linewidth}{0.4pt}
\smallskip

\noindent \textit{Instructions: Attempt all questions. Questions of Sections A, B and C carry 10, 05 and 02 marks each respectively. Write your Roll No.\ at the top immediately on receipt of this question paper.}

\smallskip
\noindent \textit{सभी प्रश्नों के उत्तर दीजिए। खण्ड अ, ब तथा स के प्रश्न क्रमशः 10, 05 एवं 02 अंकों के हैं।}

\bigskip

\section*{SECTION -- A / खण्ड -- अ}
\noindent \textit{Note: Long Answer Type Questions. Answer approximately in 250 words.}\pts{3 $\times$ 10 = 30}

\noindent \textit{दीर्घ उत्तरीय प्रश्न। लगभग 250 शब्दों में उत्तर दीजिए।}

\begin{enumerate}[label=\textbf{\arabic*.},leftmargin=*,itemsep=12pt]

  \item Trace the development of mint-towns in Ancient India.\pts{10}
  
  प्राचीन भारत में टकसाल नगरों के विकास को रेखांकित कीजिए।

  \begin{center}\textbf{OR / अथवा}\end{center}

  Write a short note on the important mints of the Post-Mauryan period.\pts{10}
  
  मौर्योत्तर काल की प्रमुख टकसालों पर एक संक्षिप्त टिप्पणी लिखिए।

  \item Describe in detail the significance of Punch-Marked coins.\pts{10}
  
  आहत सिक्कों (पंच-मार्क सिक्कों) के महत्त्व का विस्तार से वर्णन कीजिए।

  \begin{center}\textbf{OR / अथवा}\end{center}

  Describe the transition from barter to the use of metallic currency in Ancient India.\pts{10}
  
  प्राचीन भारत में वस्तु-विनिमय से धातु मुद्रा के उपयोग तक के परिवर्तन का वर्णन कीजिए।

  \item Discuss the concept of die-studies in numismatics.\pts{10}
  
  मुद्राशास्त्र में ठप्पा-अध्ययन (die-studies) की अवधारणा की चर्चा कीजिए।

  \begin{center}\textbf{OR / अथवा}\end{center}

  Highlight the technological differences and historical importance of various coin minting techniques used in Ancient India.\pts{10}
  
  प्राचीन भारत में प्रयुक्त सिक्का ढलाई की विभिन्न तकनीकों के तकनीकी अन्तर एवं ऐतिहासिक महत्त्व पर प्रकाश डालिए।

\end{enumerate}

\bigskip

\section*{SECTION -- B / खण्ड -- ब}
\noindent \textit{Note: Short Answer Type Questions. Answer approximately in 125--150 words.}\pts{4 $\times$ 5 = 20}

\noindent \textit{लघु उत्तरीय प्रश्न। लगभग 125--150 शब्दों में उत्तर दीजिए।}

\begin{enumerate}[label=\textbf{\arabic*.},leftmargin=*,start=4,itemsep=10pt]

  \item Write a short note on the coinage system of the Harappan age.\pts{5}
  
  हड़प्पा काल की विनिमय/सिक्का प्रणाली पर एक संक्षिप्त टिप्पणी लिखिए।

  \begin{center}\textbf{OR / अथवा}\end{center}

  Define obverse and reverse of coins with examples.\pts{5}
  
  सिक्कों के मुख्य भाग (Obverse) तथा पृष्ठ भाग (Reverse) को सोदाहरण परिभाषित कीजिए।

  \item Throw light on the legends on coins. Why are they important?\pts{5}
  
  सिक्कों पर उत्कीर्ण लेखों (Legends) पर प्रकाश डालिए। वे क्यों महत्त्वपूर्ण हैं?

  \begin{center}\textbf{OR / अथवा}\end{center}

  Write a short note on Takshashila as a mint centre.\pts{5}
  
  टकसाल केंद्र के रूप में तक्षशिला पर एक संक्षिप्त टिप्पणी लिखिए।

  \item Write an essay on the die-struck method of minting coins.\pts{5}
  
  सिक्का ढलाई की ठप्पा-विधि (Die-struck method) पर एक निबन्ध लिखिए।

  \begin{center}\textbf{OR / अथवा}\end{center}

  Discuss the concept of Gupta and early medieval mints in Ancient India.\pts{5}
  
  प्राचीन भारत में गुप्त एवं पूर्व मध्यकालीन टकसालों की अवधारणा की चर्चा कीजिए।

  \item Throw light on the Khokhrakot (Rohtak) mint.\pts{5}
  
  खोखराकोट (रोहतक) टकसाल पर प्रकाश डालिए।

  \begin{center}\textbf{OR / अथवा}\end{center}

  Write a short note on pre-historic exchange system.\pts{5}
  
  प्रागैतिहासिक विनिमय प्रणाली पर एक संक्षिप्त टिप्पणी लिखिए।

\end{enumerate}

\bigskip

\section*{SECTION -- C / खण्ड -- स}
\noindent \textit{Note: Very Short Answer Type Questions. Answer approximately in 15--20 words.}\pts{10 $\times$ 2 = 20}

\noindent \textit{अति लघु उत्तरीय प्रश्न। लगभग 15--20 शब्दों में उत्तर दीजिए।}

\begin{enumerate}[label=\textbf{\arabic*.},leftmargin=*,start=8,itemsep=8pt]
  \item 
  \begin{parts}
    \item What is `Nishka'?\pts{2}
    
    `निष्का' (Nishka) क्या है? (Ancient gold ornament / unit of currency)

    \item Which metal is mostly used in making Punch-marked coins?\pts{2}
    
    आहत सिक्के बनाने में किस धातु का सर्वाधिक प्रयोग किया जाता है? (Silver / चाँदी)

    \item Write the name of a book written by Upendra Thakur.\pts{2}
    
    उपेन्द्र ठाकुर द्वारा लिखित एक पुस्तक का नाम लिखिये। (Mints and Minting in Ancient India)

    \item Nalanda mint is associated with which period?\pts{2}
    
    नालंदा टकसाल किस काल से सम्बद्ध है? (Gupta / Pala period)

    \item Which is the oldest coin minting technique of Ancient India?\pts{2}
    
    प्राचीन भारत की सबसे पुरानी सिक्का निर्माण तकनीक कौन-सी है? (Punch-marked technique)

    \item What was the dynastic symbol of the Pandyas on coins?\pts{2}
    
    सिक्कों पर पाण्ड्य राजाओं का राजवंशीय चिह्न क्या था? (Fish / मत्स्य)

    \item How many symbols are generally found on imperial punch-marked coins?\pts{2}
    
    मौर्य साम्राज्यीय आहत सिक्कों पर सामान्यतः कितने चिह्न पाए जाते हैं? (Five symbols / पाँच प्रतीक)

    \item Which animal is depicted mostly on Harappan seals?\pts{2}
    
    हड़प्पा की मुहरों पर मुख्यतः किस पशु का अंकन मिलता है? (Unicorn / एकशृंगी वृषभ)

    \item Write the names of two important mints of North-Western India.\pts{2}
    
    उत्तर-पश्चिमी भारत की दो महत्त्वपूर्ण टकसालों के नाम लिखिये। (Takshashila, Pushkalavati)

    \item What do you understand by the term `Gopuchha'?\pts{2}
    
    `गोपुच्छ' (Gopuchha) शब्द से आप क्या समझते हैं?
  \end{parts}
\end{enumerate}

\end{document}
